

% = Paragraph 1 — Concluding remarks on the issue 
% = (satisfies: "one or two concluding remarks")

% Briefly restate the core problem: that Gen AI has transformed an old problem — wartime propaganda — into something qualitatively different in scale and epistemic impact. Tie back to the ethical dilemma from your introduction to give the essay a sense of closure.

The fictional ethical dilemma from the introduction section is not far-fetched given the documented cases discussed in the Issue section. Generative Artificial Intelligence has transformed the old problem of wartime propaganda into something both qualitatively and quantitatively different. The erosion of trust in visual trust is a fundamentally new issue, even if the underlying problem of propaganda is not. However, this issue is not caused by AI alone. It merely amplifies societal conditions like polarisation or a weak media literacy.


% = Paragraph 2 — Summary of the solution findings
% = (satisfies: "findings on the existing/proposed solution")

% Synthesise what your four sources collectively suggest about the state of solutions: detection lags behind generation, regulation is still catching up, and societal mechanisms are unevenly distributed globally.

According to \citet{muam2026role} Detection tools exist but they are reactive by nature, meaning they lag behind the development of increasingly improving deepfake tools. From a legal perspective, \citet{laws14060083} reveals a significant gap in International Humanitarian Law, which was not designed with AI-generated deception in mind. Existing frameworks, like the Geneva Convention, are difficult to apply in the case of deepfakes. The DFRLab reports that even well-documentented cases of deepfake-based disinformation in Ukraine did not prevent confusion among the public~\citep{osadchuk2024ai}. This shows that the effort of current societal mechanisms, like media literacy or fact-checking organisations, are not sufficient and is being outpaced by the speed at which deepfakes spread over social media.



% = Paragraph 3 — Personal answer to the core question
% = (satisfies: "is the solution suitable or  should it be a total change in approach?")

% Answer directly whether the existing solutions are sufficient or whether a fundamentally different approach is needed. Given your topic, a strong position here might be that incremental fixes are not enough and that international coordination — legal, technical, and educational — is required.

I don't believe a total change in approach is necessary, but rather that the existing solutions need to be reviewed and adjusted accordingly. Technical detection needs to be embedded into platforms at the point of distribution, and progress is being made in this direction. International Humanitarian Law should be updated to explicitly address the issue of AI-generated deception. Finally, more effort should be invested in increasing media literacy, especially in vulnerable regions such as conflict zones.

The pace of AI development will only intensify the issue of deepfakes, so coordinated action is urgently needed rather than optional.